% ===== §9 =====
\section{Optimizationism as a Symptom of Modernity}
\label{p25:sec:symptom}

\noindent
The criticism now widens from the damage optimization does within a single relation to the standing of optimizationism as a general thesis about the world. The claim of this section is that the elevation of optimization into a law of the universe is best understood not as a discovery but as a \emph{symptom}: an epistemological form taken by a particular historical condition, which the mother paper of this series named the crisis of the symbolic, and which is bound up with the habituation of a life-world to the logic of capital \parencite{wan2024grb}. The universe, on this reading, climbs no landscape; the climbing of landscapes is a historically specific practice, native to a specific organization of production, that has been projected outward onto the cosmos and taken for its law. To argue this is not to refute optimizationism by its origins, the fallacy already disowned, but to relocate it: to move it from the register of eternal truth, where it defends itself as simply how things are, into the register of historical symptom, where one may ask what condition it expresses and what that condition costs. The genealogy of the first part showed that optimizationism was constructed; this section asks what, in the present, keeps the construction standing and makes it feel like nature.

\subsection*{The commensurating operation as the native form of capital}

The decisive feature of optimizationism, isolated already in the genealogy, is its demand for commensuration: its requirement that heterogeneous goods be brought onto a single scale before a maximum can be defined. This requirement is not peculiar to the theory of optimization, and seeing where else it lives is the key to the diagnosis. It is the native operation of capital. The circulation of capital consists in nothing else than the reduction of qualitatively incomparable things, an hour of labor, a measure of grain, a song, an afternoon of care, to a common quantity in which they can be exchanged, accumulated, and made to breed more of themselves; the value-form is precisely a commensurating machine, and the historical career of capital is the spread of its commensuration into region after region of life, until things once held to lie beyond price acquire one and enter the circuit \parencite{marx1976capital}. Optimizationism is the epistemology proper to this machine. It is what universal commensuration looks like once it is raised from a practice to a theory of rationality and thence to a picture of the world: the settled conviction that behind the apparent heterogeneity of goods there always lies a scale on which they compare and a maximum toward which reason should climb. The conviction feels self-evident to the modern mind not because it is a self-evident truth but because the mind has been formed within a life-world that performs the commensurating operation on everything, ceaselessly, until the performability of the operation comes to seem a property of the goods rather than an imposition laid upon them. This is the critique of instrumental reason that the Frankfurt tradition pressed: that reason, reduced to the calculation of means and the commensuration of ends, mistakes its own historically produced form for reason as such and forgets the ends it can no longer reason about \parencite{horkheimer2004eclipse,adornohorkheimer2002}.

\subsection*{The symbolic crisis and the flight to the scale}

The mother paper diagnosed the present as a crisis of the symbolic, a hollowing of the shared meanings by which a world is held together, and optimizationism can be placed within that diagnosis as one of its characteristic responses \parencite{wan2024grb}. The placement explains something the genealogy alone could not: why the constructed picture is not merely available but attractive, why people reach for it. When the thick symbolic orders that once told a community what things were worth, and told it in qualitative and incommensurable terms, lose their hold, a vacuum opens where the sense of worth had been, and into that vacuum the single quantitative scale flows as a replacement. A scale asks for no shared meaning. It needs no agreement about what a thing is for or what makes a life go well, only a number; and so it promises orientation to a disenchanted world that can no longer supply itself with meanings, a way to know what to do, maximize the measure, that survives the collapse of every thick account of the goods. Optimizationism is in this sense a flight: a flight from the difficulty of incommensurable value, which requires judgment and yields no algorithm, into the false comfort of a scale, which requires only computation and spares one the judgment. That the flight should present itself as hard-headed rationality, and the refusal of it as soft sentiment, is itself part of the symptom, for the condition valorizes the very operation that expresses it and recasts fidelity to the real structure of value, the refusal to commensurate what does not compare, as a failure of nerve to be overcome.

\subsection*{The projection onto the cosmos}

The last movement of the symptom is the projection of the historically specific operation onto the universe as its law, and it is this movement that the paper's title refuses. Because the commensurating, maximizing operation is performed everywhere within a certain form of life, it comes to seem that the operation is not being performed but merely recognized, that the landscapes were always there and reason only climbs them. One meets this projection in its popular cosmic form, in which everything from a molecule seeking its lowest energy to a market clearing to an organism maximizing its fitness is said to be ``really'' optimizing some quantity, so that the human who calculates is at last only doing consciously what all of nature does blindly, and optimization stands revealed as the deep grammar of being. The paper's denial is aimed exactly here, and it can be put precisely. What is actually observed is not a universe that climbs landscapes but a form of life that climbs them and then, carrying its scale with it, finds the operation everywhere it lays the scale down. That certain physical systems admit description as extremizing a quantity is a fact about the reach of a mathematical formalism, not evidence that nature is engaged in maximization in the sense at issue, the sense in which there is an objective owned by a subject and a best state to be sought; the formalism's success has been read back into the world as the world's own procedure. To say that optimization is a law of the universe is to mistake the shadow the commensurating life-world casts over whatever it examines for a feature of what is examined. The practice is real and has a date. The law is a projection. And intimate relation, whose goods are relational and incommensurable and whose subjects are constituted in the relating, is the region where the projection is most plainly false, because it is the region in which the commensurating operation, once applied, most visibly destroys the very thing it was applied to.
